\chapter{Notificaciones push reales con Firebase Cloud Messaging}
\label{ch:fcm}

\section{Objetivo}

Reemplazar la simulacion de notificaciones push de la Sesion~\ref{ch:sesion5} por un sistema real de Firebase Cloud Messaging (FCM) que:

\begin{itemize}[leftmargin=*]
  \item Registre cada dispositivo del usuario contra el backend.
  \item Envie notificaciones reales cuando se publique en un canal suscrito.
  \item Muestre notificaciones en los tres estados de la app: foreground, background y quit.
  \item Permita navegar al canal correspondiente al tocar la notificacion.
\end{itemize}

\section{Contexto}

Al cierre de la Sesion~\ref{ch:sesion5}, el \texttt{PostNotificationConsumer} enviaba notificaciones simuladas con logs del tipo \texttt{[PUSH-SIMULADO]}. El objetivo final del protocolo (Fase~6) exige notificaciones push reales. Se dividio el trabajo en cuatro sub-fases:

\begin{enumerate}[leftmargin=*]
  \item \textbf{F1:} Configuracion del proyecto en Firebase Console.
  \item \textbf{F2:} Backend con Firebase Admin SDK y entidad \texttt{DeviceToken}.
  \item \textbf{F3:} Configuracion Android + Flutter con \texttt{firebase\_messaging}.
  \item \textbf{F4:} Manejo de notificaciones en foreground, background y quit.
\end{enumerate}

\section{Desarrollo}

\subsection{F1: Configuracion en Firebase Console}

Se creo un proyecto en Firebase Console llamado \texttt{notiuam-4ee15} y se registro una app Android con el paquete \texttt{mx.uam.notiuam\_mobile}. De la consola se descargaron dos archivos:

\begin{itemize}[leftmargin=*]
  \item \texttt{google-services.json}: se coloco en \texttt{android/app/}.
  \item \texttt{service-account.json}: se coloco en \texttt{backend/src/main/resources/firebase/}.
\end{itemize}

Ambos archivos se agregaron a \texttt{.gitignore} porque contienen credenciales privadas.

\subsection{F2: Backend con Firebase Admin SDK}

\subsubsection{Dependencia}

Se agrego al \texttt{pom.xml}:

\begin{lstlisting}[language=XML]
<dependency>
    <groupId>com.google.firebase</groupId>
    <artifactId>firebase-admin</artifactId>
    <version>9.4.1</version>
</dependency>
\end{lstlisting}

\subsubsection{Migracion V3}

\begin{lstlisting}[language=SQL]
CREATE TABLE device_tokens (
    id          UUID PRIMARY KEY DEFAULT gen_random_uuid(),
    user_id     UUID         NOT NULL REFERENCES users(id) ON DELETE CASCADE,
    token       VARCHAR(255) NOT NULL UNIQUE,
    platform    VARCHAR(20)  NOT NULL CHECK (platform IN ('ANDROID', 'IOS', 'WEB')),
    created_at  TIMESTAMPTZ  NOT NULL DEFAULT NOW(),
    updated_at  TIMESTAMPTZ  NOT NULL DEFAULT NOW()
);

CREATE INDEX idx_device_tokens_user ON device_tokens(user_id);
\end{lstlisting}

\subsubsection{Inicializacion del SDK}

\texttt{FirebaseConfig} carga el service account al arrancar la aplicacion y lo deja listo para usar.

\begin{lstlisting}[language=Java]
@Configuration
public class FirebaseConfig {
    @PostConstruct
    public void initialize() {
        try {
            if (!FirebaseApp.getApps().isEmpty()) return;

            ClassPathResource resource = new ClassPathResource(
                "firebase/service-account.json");
            if (!resource.exists()) {
                log.warn("Sin service-account.json. Push deshabilitado.");
                return;
            }

            try (InputStream is = resource.getInputStream()) {
                FirebaseOptions options = FirebaseOptions.builder()
                    .setCredentials(GoogleCredentials.fromStream(is))
                    .build();
                FirebaseApp.initializeApp(options);
                log.info("Firebase Admin SDK inicializado");
            }
        } catch (IOException e) {
            log.error("Error inicializando Firebase: {}", e.getMessage());
        }
    }
}
\end{lstlisting}

\subsubsection{Servicio de envio}

\texttt{FcmService} encapsula la logica de envio:

\begin{lstlisting}[language=Java]
public boolean sendToToken(String deviceToken, String title,
                           String body, Map<String, String> data) {
    if (!isEnabled()) return false;

    try {
        Message message = Message.builder()
            .setToken(deviceToken)
            .setNotification(Notification.builder()
                .setTitle(title)
                .setBody(body)
                .build())
            .putAllData(data != null ? data : Map.of())
            .build();

        String response = FirebaseMessaging.getInstance().send(message);
        return true;
    } catch (FirebaseMessagingException e) {
        log.warn("Error enviando push: {} ({})",
            e.getMessage(), e.getMessagingErrorCode());
        return false;
    }
}
\end{lstlisting}

\subsubsection{Endpoints}

\begin{lstlisting}[language=bash]
POST   /api/me/device              (registrar token)
DELETE /api/me/device?token=...    (eliminar token)
\end{lstlisting}

\subsubsection{Modificacion a \texttt{PostNotificationConsumer}}

\begin{aviso}[Cambio a un archivo previo]
El consumidor de la Sesion~\ref{ch:sesion5} se modifico para enviar notificaciones reales en lugar de solo loguear. Se inyectaron \texttt{FcmService} y \texttt{DeviceTokenService} y se anadio el envio por token:

\begin{lstlisting}[language=Java]
List<String> tokens = deviceTokens.tokensOf(sub.getUserId());
for (String token : tokens) {
    boolean ok = fcm.sendToToken(
        token,
        event.channelName(),
        truncate(event.content()),
        Map.of(
            "channelId", event.channelId().toString(),
            "postId", event.postId().toString()
        )
    );
    if (ok) log.info("[PUSH-FCM] -> user={}", sub.getUserId());
}
\end{lstlisting}
\end{aviso}

\subsection{F3: Configuracion Android y Flutter}

\subsubsection{Gradle}

Se agrego el plugin de Google Services en \texttt{settings.gradle.kts} y \texttt{app/build.gradle.kts}.

\begin{lstlisting}[language=Java]
// settings.gradle.kts
plugins {
    id("dev.flutter.flutter-plugin-loader") version "1.0.0"
    id("com.android.application") version "9.1.0" apply false
    id("org.jetbrains.kotlin.android") version "2.4.0" apply false
    id("com.google.gms.google-services") version "4.4.2" apply false
}

// app/build.gradle.kts
plugins {
    id("com.android.application")
    id("dev.flutter.flutter-gradle-plugin")
    id("com.google.gms.google-services")
}

android {
    compileOptions {
        isCoreLibraryDesugaringEnabled = true
        sourceCompatibility = JavaVersion.VERSION_17
        targetCompatibility = JavaVersion.VERSION_17
    }
}

dependencies {
    coreLibraryDesugaring("com.android.tools:desugar_jdk_libs:2.1.4")
}
\end{lstlisting}

\subsubsection{Permisos en AndroidManifest}

\begin{lstlisting}[language=XML]
<uses-permission android:name="android.permission.INTERNET"/>
<uses-permission android:name="android.permission.POST_NOTIFICATIONS"/>
<uses-permission android:name="android.permission.VIBRATE"/>
\end{lstlisting}

Y dentro de \texttt{<application>}:

\begin{lstlisting}[language=XML]
<meta-data
    android:name="com.google.firebase.messaging.default_notification_channel_id"
    android:value="high_importance_channel" />
\end{lstlisting}

\subsubsection{Dependencias Flutter}

Se agregaron a \texttt{pubspec.yaml}:

\begin{lstlisting}[language=YAML]
firebase_core: ^3.6.0
firebase_messaging: ^15.1.3
flutter_local_notifications: ^18.0.1
\end{lstlisting}

\subsubsection{\texttt{FcmService}}

Servicio encargado de pedir permisos, obtener el token y registrarlo en el backend:

\begin{lstlisting}[language=Java]
Future<void> initialize() async {
    if (!Platform.isAndroid && !Platform.isIOS) return;

    final messaging = FirebaseMessaging.instance;
    final settings = await messaging.requestPermission(
        alert: true, badge: true, sound: true);

    if (settings.authorizationStatus == AuthorizationStatus.denied) return;

    final token = await messaging.getToken();
    if (token == null) return;

    await _registerWithBackend(token);
    messaging.onTokenRefresh.listen(_registerWithBackend);
    FirebaseMessaging.onMessage.listen(_onForegroundMessage);
    FirebaseMessaging.onMessageOpenedApp.listen(_onMessageOpened);

    final initial = await messaging.getInitialMessage();
    if (initial != null) _onMessageOpened(initial);
}

Future<void> _registerWithBackend(String token) async {
    final platform = Platform.isAndroid ? 'ANDROID' : 'IOS';
    await _api.dio.post('/api/me/device', data: {
        'token': token, 'platform': platform,
    });
}
\end{lstlisting}

\subsubsection{Integracion con \texttt{AuthProvider}}

El token se registra automaticamente al hacer login y se elimina al cerrar sesion:

\begin{lstlisting}[language=Java]
Future<void> _handleAuthResponse(dynamic data) async {
    // ... guardar tokens y usuario ...
    await _fcm.initialize();   // registrar dispositivo
}

Future<void> logout() async {
    await _fcm.unregister();   // eliminar token
    // ...
}
\end{lstlisting}

\subsection{F4: Manejo de notificaciones}

\subsubsection{Los tres estados}

FCM distingue tres estados de la aplicacion:

\begin{table}[htbp]
\centering
\begin{tabular}{l l}
\toprule
Estado & Callback \\
\midrule
Foreground (app abierta) & \texttt{FirebaseMessaging.onMessage} \\
Background (app minimizada) & \texttt{FirebaseMessaging.onMessageOpenedApp} \\
Quit (app cerrada) & \texttt{messaging.getInitialMessage()} \\
\bottomrule
\end{tabular}
\caption{Callbacks de FCM segun el estado de la aplicacion}
\end{table}

Los tres se manejan en \texttt{FcmService.initialize()} y todos disparan el mismo callback \texttt{onNotificationTap} que navega al canal.

\subsubsection{Background handler}

Para que el sistema operativo pueda procesar mensajes cuando la app esta cerrada, se registra un handler global:

\begin{lstlisting}[language=Java]
@pragma('vm:entry-point')
Future<void> firebaseMessagingBackgroundHandler(RemoteMessage message) async {
    await Firebase.initializeApp();
    debugPrint('FCM background: ${message.notification?.title}');
}
\end{lstlisting}

Se registra en \texttt{main()}:

\begin{lstlisting}[language=Java]
FirebaseMessaging.onBackgroundMessage(firebaseMessagingBackgroundHandler);
\end{lstlisting}

\subsubsection{Notificaciones locales}

Para mostrar notificaciones cuando la app esta en foreground (FCM por defecto no lo hace), se usa \texttt{flutter\_local\_notifications}:

\begin{lstlisting}[language=Java]
void _onForegroundMessage(RemoteMessage message) {
    final n = message.notification;
    if (n != null) {
        LocalNotifications.show(
            title: n.title ?? 'NotiUAM',
            body: n.body ?? '',
        );
    }
}
\end{lstlisting}

El canal \texttt{high\_importance\_channel} se crea al arrancar la app:

\begin{lstlisting}[language=Java]
const channel = AndroidNotificationChannel(
    'high_importance_channel',
    'Notificaciones',
    description: 'Avisos y alertas de NotiUAM',
    importance: Importance.high,
);
await _plugin
    .resolvePlatformSpecificImplementation<AndroidFlutterLocalNotificationsPlugin>()
    ?.createNotificationChannel(channel);
\end{lstlisting}

\subsubsection{Navegacion al canal}

Para que el toque en la notificacion abra el canal correcto, se usa un \texttt{GlobalKey<NavigatorState>}:

\begin{lstlisting}[language=Java]
final navigatorKey = GlobalKey<NavigatorState>();

void _wireFcm(FcmService fcm) {
    fcm.onNotificationTap = (channelId, postId) {
        navigatorKey.currentState?.push(MaterialPageRoute(
            builder: (_) => ChannelDetailScreen(channelId: channelId),
        ));
    };
}
\end{lstlisting}

\subsection{Cambios a archivos previos}

\begin{itemize}[leftmargin=*]
  \item \texttt{pom.xml}: se agrego la dependencia \texttt{firebase-admin}.
  \item \texttt{PostNotificationConsumer.java}: se reemplazo la simulacion por el envio real.
  \item \texttt{pubspec.yaml}: se agregaron las tres dependencias de Firebase.
  \item \texttt{main.dart}: se inicializa Firebase, se registra el background handler y se inyecta \texttt{FcmService}.
  \item \texttt{auth\_provider.dart}: \texttt{\_fcm.initialize()} al login y \texttt{\_fcm.unregister()} al logout.
  \item \texttt{settings.gradle.kts} y \texttt{app/build.gradle.kts}: plugins de Firebase.
  \item \texttt{AndroidManifest.xml}: permisos y meta-data del canal.
\end{itemize}

\subsection{Problemas encontrados y solucion}

\subsubsection{Falta de \texttt{uses-material-design}}

Al inicio los iconos de Material Design se veian como glifos de fuentes CJK. Se corrigio en la Fase de rediseno agregando \texttt{uses-material-design: true} al \texttt{pubspec.yaml}.

\subsubsection{Core library desugaring}

\texttt{flutter\_local\_notifications} requiere desugaring de librerias de Java 8:

\begin{lstlisting}[language=bash]
Dependency ':flutter_local_notifications' requires core library
desugaring to be enabled for :app.
\end{lstlisting}

\textbf{Solucion:} se agrego \texttt{isCoreLibraryDesugaringEnabled = true} y la dependencia \texttt{coreLibraryDesugaring(...)} en \texttt{app/build.gradle.kts}.

\subsubsection{Dependencias de Firebase no declaradas}

Inicialmente \texttt{flutter pub get} no descargaba \texttt{firebase\_core} ni \texttt{firebase\_messaging}. Faltaba declararlas explicitamente en \texttt{pubspec.yaml}.

\subsubsection{Background handler no registrado}

FCM advertia:

\begin{lstlisting}[language=bash]
A background message could not be handled in Dart as no
onBackgroundMessage handler has been registered.
\end{lstlisting}

\textbf{Solucion:} crear \texttt{fcm\_background.dart} con el handler y registrarlo en \texttt{main()} con \texttt{FirebaseMessaging.onBackgroundMessage()}.

\subsubsection{Canal de notificaciones no creado}

FCM advertia que el canal declarado en el manifest no existia en la app. Se resolvio creandolo explicitamente al arrancar con \texttt{flutter\_local\_notifications}.

\section{Verificacion}

\subsection{Registro del token}

Al hacer login en el emulador Android aparecen los logs:

\begin{lstlisting}[language=bash]
I/flutter: FCM token: f4Pr5kGlQw-IsWljYQ6E...
I/flutter: FCM: token registrado en el backend
\end{lstlisting}

Y en la base de datos:

\begin{lstlisting}[language=SQL]
SELECT u.email, d.platform, LEFT(d.token, 25) FROM device_tokens d
JOIN users u ON u.id = d.user_id;

    email      | platform |       token_prefix
---------------+----------+---------------------------
 aaron@uam.mx  | ANDROID  | f4Pr5kGlQw-IsWljYQ6EdM:AP
\end{lstlisting}

\subsection{App en foreground}

Con la app abierta, cuando otro usuario publica en un canal suscrito:

\begin{lstlisting}[language=bash]
I/flutter: FCM foreground: <nombre del canal> - <contenido>
\end{lstlisting}

Y aparece una notificacion local visible en el emulador.

\subsection{App en background}

Con la app minimizada, la notificacion aparece en la barra del sistema Android. Al tocarla:

\begin{lstlisting}[language=bash]
I/flutter: FCM tap: channelId=0c1c26c1-... postId=b0738e37-...
\end{lstlisting}

Y la app navega automaticamente al canal.

\subsection{App cerrada (quit)}

Con la app cerrada completamente, al publicar aparece la notificacion. Al tocarla, la app arranca y despues de 500ms navega al canal.

\subsection{Logs del backend}

Despues de cada publicacion, los logs muestran:

\begin{lstlisting}[language=bash]
[PUSH-FCM] -> token=f4Pr5kGlQw... user=da4e045f...
Notificaciones procesadas: 1 de N suscriptores
\end{lstlisting}

El contador \texttt{1 de N} excluye al autor del mensaje.

\section{Resultado}

\begin{exito}[Estado al cierre]
\begin{itemize}[leftmargin=*]
  \item Sistema de push real con FCM end-to-end.
  \item Token registrado automaticamente al login y eliminado al logout.
  \item Notificaciones en los tres estados de la app.
  \item Tap en notificacion navega al canal correspondiente.
  \item Canal \texttt{high\_importance\_channel} creado en el arranque.
  \item Reemplazo completo del \texttt{[PUSH-SIMULADO]} por \texttt{[PUSH-FCM]}.
\end{itemize}
\end{exito}
